% !TEX root = ../report.tex
\structchapter{ЗАКЛЮЧЕНИЕ}
Цель работы достигнута: выполнен сравнительный анализ ресурсной эффективности четырёх алгоритмов поиска в массиве.

В ходе работы решены поставленные задачи:
\begin{itemize}
\item описаны входные данные и результат поиска;
\item разработаны линейный, два нерекурсивных бинарных и рекурсивный бинарный алгоритмы;
\item построены схемы основных функций;
\item сформулированы требования к программе;
\item выбраны язык программирования и инструменты сборки;
\item реализованы алгоритмы и программа для их сравнения;
\item выполнено функциональное тестирование и проверка счётчиков;
\item оценена дополнительная память, включая цепочку рекурсивных вызовов;
\item получены оценки числа сравнений в лучшем и худшем случаях;
\item результаты теоретического анализа и эксперимента представлены в таблицах;
\item сформулированы выводы о ресурсных затратах и ограничениях алгоритмов.
\end{itemize}

Линейный поиск не требует упорядоченности и использует постоянную дополнительную память, но в худшем случае выполняет $N$ сравнений. Все рассмотренные бинарные алгоритмы требуют упорядоченного массива и имеют логарифмическое число сравнений в худшем случае. Классический и рекурсивный варианты могут завершить поиск сразу после обнаружения равенства. Алгоритм без ветвления сначала определяет позицию первого элемента, не меньшего искомого, а затем проверяет равенство.

Нерекурсивные реализации используют $\Theta(1)$ дополнительной памяти, рекурсивная --- $\Theta(\log N)$ в худшем случае. Эксперимент подтвердил выведенные формулы на исследованных размерах. Время выполнения программ в этой работе не измерялось.
